\subsection{Determinant Condition and When the Method Is Useful}

The inverse method applies exactly when
\[
\det A\neq0.
\]
If $\det A=0$, no inverse exists. The system may then have no solution or infinitely many solutions, but it cannot be solved by multiplying by $A^{-1}$.

\begin{examplebox}
For
\[
\begin{cases}
x+y=2,\\
2x+2y=4,
\end{cases}
\]
the coefficient determinant is zero, yet every point on $x+y=2$ is a solution. Failure of invertibility means loss of uniqueness, not automatic inconsistency.
\end{examplebox}

\begin{interpretationbox}[When inversion is efficient]
The method is especially useful when $A^{-1}$ is already known, when several right-hand sides share the same coefficient matrix, or when the system is being studied explicitly as a matrix transformation. Direct elimination is often cheaper for a single routine system.
\end{interpretationbox}

\begin{summarybox}[Inverse method checklist]
Write $Ax=b$, check that $A$ is square and invertible, multiply on the left by $A^{-1}$, and verify the resulting vector in the original equations.
\end{summarybox}
